\section{Finite projected-factor catalogues and a large-field tensor obstruction}
\label{sec:tensor-factor-catalogue}

\paragraph{Scope and status.}
D-TD01--03 concern an actual factorwise tensor decomposition.
They do not transfer the twelve-mark condition to a projected factor.
The finite catalogue statement is conditional only on bounds for the
projective factor images; the large-field statement below is an
unconditional exclusion under its explicit normal-core hypotheses.
These are new \texttt{/root/imprimitive} derivations, discussed with
root. Complete author and separate root whole-proof readings pass;
root also read the full source/status record, checker, output and
invoked primary DHP and Milne passages. Exact boundaries are recorded
in the retained review. No global exhaustion of all factor types is asserted.

Let \(J\subset\operatorname{GL}(V)\) be a finite faithful irreducible
matrix group of degree \(e\ge2\).  Suppose an actual connected sphere
cover has a genuine generating product-one inertia tuple
\(g_1,\ldots,g_r\in J\), satisfying
\[
 r\ge12,\qquad
 \sum_{i=1}^r\operatorname{rank}(1-g_i)=2e+6 .
 \tag{TD.1}
\]
For the tensor catalogue assume that \(J\) preserves factorwise
\[
                    V=W_1\otimes\cdots\otimes W_t,
             \qquad a_j=\dim W_j\ge2 .
                    \tag{TD.2}
\]
The resulting projective factor homomorphisms are denoted
\(J\longrightarrow P_j\subset\operatorname{PGL}(W_j)\).
They are actual homomorphisms: the factors in a tensor decomposition
of an invertible operator are unique up to scalars.
Every projective factor action is irreducible, since an invariant
proper subspace in one factor would give one in \(V\).

\subsection{D-TD01: bounded projective factors bound the whole inducing image}

\begin{theorem}[D-TD01]
Assume \(a_j\le A\) and \(|P_j|\le B\) for every \(j\), where
\(A,B\ge2\) are integers.  Define
\[
 R=5A,\qquad
 N=B^{B^2+R+1},\qquad
 M=B^N,\qquad Q=M^3 .
 \tag{TD.3}
\]
Then
\[
            r\le R,\qquad
       |J/(J\cap\mathbf C^\times I)|\le M,\qquad |J|\le Q .
       \tag{TD.4}
\]
There is no bound on \(t\) among the hypotheses.
Consequently the actual faithful inducing images, their irreducible
modules, and generating branch tuples satisfying these bounds lie in
an effectively enumerable finite list.
\end{theorem}

\begin{proof}
For a finite-order operator \(T\), write
\(\nu(T)=\dim T-\max_z\dim\ker(T-z)\).
Tensor spectral convolution gives
\[
 \frac{\operatorname{rank}(1-g_i)}e
             \ge \max_j\frac{\nu(W_j(g_i))}{a_j}.
 \tag{TD.5}
\]
Representatives of the projective operators may be chosen arbitrarily;
their scalar phases do not affect \(\nu\).
If some projected operator is nonidentity, it has \(\nu\ge1\), so
the right side is at least \(1/A\).
If all are projectively trivial, \(g_i\) is scalar; because the tuple
is genuine, its actual moving rank is \(e\), again at least \(e/A\).
Thus (TD.1) gives
\[
                  r\le A(2+6/e)\le5A=R .
                  \tag{TD.6}
\]

Pad the tuple by identities to length \(R\), and let \(F_R\) be the
free group on these labelled generators.  For a group of order
\(b\le B\), fix a labelling by \(\{1,\ldots,b\}\).
There are at most \(b^{b^2}\) binary operation tables and at most
\(b^R\) assignments of the \(R\) generators.  These deliberately
overcount all group structures and all homomorphisms to them.
In particular there are at most
\[
          \sum_{b=1}^B b^{b^2+R}
                     \le B^{B^2+R+1}=N
                     \tag{TD.7}
\]
distinct kernels among all such homomorphisms.
Each actual \(F_R\longrightarrow P_j\) belongs to this list.
Repeating a kernel has no effect on their intersection.
Select one representative for every distinct kernel occurring in
(TD.2); at most \(N\) representatives are needed.

An invertible tensor product is scalar only if each factor is scalar.
For example, conjugating \(X\otimes I\otimes\cdots\otimes I\)
by a scalar tensor product shows that its first factor commutes
with every \(X\), and therefore is scalar; repeat for all factors.
It follows that the kernel of the combined map to the \(P_j\)
is exactly the scalar kernel in \(J\).  Hence, putting
\(s=|J/(J\cap\mathbf C^\times I)|\), the selected quotient maps give
\[
                              s\le B^N=M .
                              \tag{TD.8}
\]
This bounds the order, not a claimed divisibility by \(M\).

We next bound the scalar subgroup without assuming any field.
Irreducibility implies that the span of \(J\) is
\(\operatorname{End}_{\mathbf C}(V)\), by the group-algebra
decomposition \cite[\S6.2, Proposition10, p.~48]{OG:Serre1977}.
Choose one matrix representative for each projective element.
Every matrix of \(J\) is a scalar multiple of one of these \(s\)
representatives.  Therefore
\[
                                  e^2\le s.
                                  \tag{TD.9}
\]

At most one inertia has no eigenvalue \(1\).  Otherwise two entries
cost \(2e\), and the other \(r-2\ge10\) genuine entries each cost
at least one, contradicting (TD.1).
For every other entry, \(g_i^s\) is scalar by Lagrange's theorem
in the projective image.  Since \(g_i\) has eigenvalue \(1\),
that scalar is one.  Thus \(\det(g_i)^s=1\).
If there is an exceptional entry, the determinant of the
product-one relation gives the same conclusion for it.
The tuple generates \(J\), so \(\det(J)\subset\mu_s\).
For \(zI\in J\), this implies \(z^{es}=1\), whence
\[
 |J\cap\mathbf C^\times I|\le es,\qquad
 |J|\le es^2\le s^{5/2}\le M^3 .
 \tag{TD.10}
\]
This proves (TD.4).

For effective enumeration, enumerate the finitely many multiplication
tables of order at most \(Q\), test the group axioms, and construct
their irreducible representations over a cyclotomic splitting field.
The existence of such a field is the known Serre theorem already
used in the finite-rank algorithm
\cite[\S12.3, Theorem24, p.~94]{OG:Serre1977}.
Finite semisimple-algebra computations, or the exact algebraic
enumeration in D-NR, give an effective construction; no undecidable
choice is required.  Enumerate tuples of length at most \(R\)
and test generation and product one in the tables.
It is harmless to retain a superset instead of deciding every
tensor-decomposition condition.  This finite list contains every
source under the stated hypotheses.
\end{proof}

\paragraph{Order-only version.}
If only \(|P_j|\le B\) is given, irreducibility still implies
\(a_j^2\le |P_j|\).  Indeed the inverse image of \(P_j\) in
\(\operatorname{SL}(W_j)\) is finite, its representation is
irreducible, and its matrix span is generated by one representative
for each projective element, exactly as in (TD.9).
Thus one may take \(A=B\) in (TD.3), or the sharper integer
\(A=\lfloor\sqrt B\rfloor\) whenever the hypotheses are nonvacuous.
A bound on projective factor orders alone therefore suffices.

\paragraph{Scope of the finite decision.}
The list is one of inducing images and their actual tuple data.
Within the precise projector and period-map input scope of D-NR,
its exact generic-rank procedure is available for each entry.
No claim here decides whether a proper special subfamily has an
original-\(K\) realization or a prescribed polarization descent.
D-OG shows that arbitrary original overgroups can remain unbounded.
The necessary determinant filter HD-PD02 applies to the order
ceiling \(Q\), but is not promoted into a sufficient descent test.
In a primitive \(J\), each projective factor in (TD.2) is primitive:
an imprimitivity decomposition in one factor tensors to an
imprimitivity decomposition of \(V\).
Thus any independently proved projective-order bounds for such
factors may be inserted into (TD.3).

\subsection{D-TD02: a projected orbifold bound replacing twelve marks}

\begin{lemma}[D-TD02]
Suppose an actual quotient \(P\) of the source group is a finite
nonsolvable group with \(|P|>60\).  Suppose
\(V=W\otimes U\), where \(a=\dim W\), and the action on \(W\)
is a projective representation of \(P\).
Assume that every nonidentity \(x\in P\), in a unitary scalar
normalization, satisfies
\[
                         |\operatorname{tr}W(x)|/a\le c<1 .
                         \tag{TD.11}
\]
Then every actual sphere source with moving sum \(2e+6\) satisfies
\[
                               e(1-85c)\le252 .
                               \tag{TD.12}
\]
In particular \(c\le1/170\) forces \(e\le504\).
This lemma needs no lower bound on the number of genuine projected
marks, and places no restriction on \(U\).
\end{lemma}

\begin{proof}
Remove identity images from the projected generating tuple, and
write \(o_1,\ldots,o_k\ge2\) for the exact orders of the remaining
elements.  The associated actual connected \(P\)-Galois cover
\(Z\longrightarrow\mathbf P^1\) has
\[
 2g(Z)-2=|P|\left(-2+\sum_{i=1}^k(1-1/o_i)\right).
 \tag{TD.13}
\]
We check that \(g(Z)\ge2\).  If \(k\le2\), product one makes
the generated group cyclic, contrary to the hypothesis.
If \(g(Z)=0\), then \(k=3\), since every summand is at least
\(1/2\).  Ordering the three orders, the strict inequality
\(\sum1/o_i>1\) gives exactly
\((2,2,m),(2,3,3),(2,3,4),(2,3,5)\).
In the first case two involutions generate a dihedral quotient,
which is solvable.  In the other cases (TD.13) gives respectively
\(|P|=12,24,60\).  All are excluded.
If \(g(Z)=1\), choose an origin on \(Z\).  Its automorphism group
is the translation group semidirect its origin-fixing group;
over \(\mathbf C\) the latter is cyclic of order \(2,4\), or \(6\).
These exact known statements, with proofs, are
\cite[ChapterIV, Proposition7.13 and Remark7.18, pp.~140--141]{TD:MilneEC2021}.
Consequently every finite subgroup is abelian-by-cyclic and solvable,
again a contradiction.

Put \(\Delta=\sum_i(1-1/o_i)-2>0\).
The elementary hyperbolic order estimate is
\[
                         \Delta\ge1/42 .
                         \tag{TD.14}
\]
For completeness, if \(k\ge5\) it is at least \(1/2\).
For \(k=4\), the all-two case has zero excess and every positive
case has excess at least \(1/6\).
For \(k=3\), put \(u\le v\le w\).
If \(u\ge4\) the excess is at least \(1/4\).
If \(u=3\), the smallest positive case is \((3,3,4)\),
with excess \(1/12\).
If \(u=2\), the cases \(v=2\) are negative; for \(v=3\),
positivity requires \(w\ge7\) and the minimum is \(1/42\);
for \(v=4\) it is \(1/20\), and for \(v\ge5\) it is at least
\(1/10\).  This proves (TD.14) for every possible number of marks.

For an element \(x\in P\) of order \(o\), choose a scalar
normalization of its unitary representative \(T\) with \(T^o=1\).
Every power \(T^j\), \(1\le j<o\), is projectively nonidentity.
Fourier inversion and (TD.11) therefore give
\[
 \frac{\nu(T)}a
 \ge1-\frac{1+(o-1)c}{o}
                    =(1-c)(1-1/o).
 \tag{TD.15}
\]
Tensor convolution cannot decrease this normalized lower bound
on actual moving rank, whatever the partner operator and scalar.
Summing and using (TD.14) yields
\[
             2+\frac6e\ge(1-c)\frac{85}{42}.
             \tag{TD.16}
\]
Rearrangement gives (TD.12).  The specialization \(c\le1/170\)
has \(1-85c\ge1/2\), proving \(e\le504\).
\end{proof}

\paragraph{All-power precision.}
Condition (TD.11) is an absolute character bound for the complete
actual projective image, hence for every nontrivial power.
A support estimate only on a normal subgroup, or a commutator
support estimate with no trace control, cannot be substituted for it.

\subsection{D-TD03: full normalizers of large-field Lie cores}

\begin{lemma}[Normalizer trace averaging]
Let \(N\subset\operatorname{GL}(W)\) be finite, perfect, faithful,
and irreducible, with \(a=\dim W>1\), and put \(S=N/Z(N)\).
Assume the actual core character has
\[
                  |\chi_W(x)|/a\le c<1\qquad(x\notin Z(N)),
                  \tag{TD.17}
\]
and every nontrivial ordinary irreducible representation of \(S\)
has degree at least \(b\).
Then every nonscalar unitary operator \(T\) normalizing \(N\)
satisfies
\[
 \frac{|\operatorname{tr}T|^2}{a^2}
                   \le c+\frac{1-c}{b+1}.
                   \tag{TD.18}
\]
The operator may induce an outer automorphism or permute central
factors of \(N\).
\end{lemma}

\begin{proof}
Let \(\alpha\) be the induced automorphism of \(S\).
If \(\alpha=1\), then
\(TxT^{-1}=\lambda(x)x\) for all \(x\in N\), with
\(\lambda:N\longrightarrow\mathbf C^\times\) a homomorphism.
Perfectness gives \(\lambda=1\), and Schur's lemma makes \(T\)
scalar.  Thus \(\alpha\ne1\).
Its fixed subgroup \(H=C_S(\alpha)\) is proper.
The coset permutation module \(\mathbf C[S/H]\) contains the
constants once and has a nonzero complementary representation.
Its latter dimension is \([S:H]-1\), so
\[
                 [S:H]\ge b+1,\qquad |H|/|S|\le1/(b+1).
                 \tag{TD.19}
\]
For \(x\in N\), the commutator \(TxT^{-1}x^{-1}\) is central
exactly when \(xZ(N)\in H\).
Schur averaging gives the exact identity
\[
 \frac1{|N|}\sum_{x\in N}
       \frac{\operatorname{tr}(TxT^{-1}x^{-1})}{a}
                         =\frac{|\operatorname{tr}T|^2}{a^2}.
                         \tag{TD.20}
\]
Indeed the average of \(xT^{-1}x^{-1}\) is
\(\operatorname{tr}(T^{-1})I/a\).
The central commutators have absolute character ratio one,
including their possible nontrivial scalar phases.
Every other commutator has ratio at most \(c\).
The triangle inequality, (TD.19), and (TD.20) prove (TD.18).
\end{proof}

\begin{theorem}[D-TD03]
Suppose an actual factor \(W\) of a sphere-source coefficient
module \(V=W\otimes U\) has a normal irreducibly acting central product
\[
       N=E_1\cdots E_t\subset\operatorname{GL}(W),
       \qquad S_i=E_i/Z(E_i),
       \tag{TD.21}
\]
where each \(E_i\) is a finite quasisimple group of Lie type with
nonalternating simple quotient, in the defining-field convention of
Duyan--Halasi--Podoski Theorem4.1(2).
Assume every defining field size satisfies \(q_i\ge10^{12}\).
The actual factor image may be any finite group in the full linear
normalizer of \(N\), including compatible outer automorphisms,
factor permutations and arbitrary scalar extensions.
Then no such sphere source has moving sum \(2e+6\).
There is no factorwise tensor-indecomposability hypothesis, no
transitivity requirement on the \(E_i\), no original-field restriction,
and no projected twelve-mark premise.
\end{theorem}

\begin{proof}
The precise published character estimate is
\[
 {\rm mr}(E_i)\le
       \max\{1/(\sqrt{q_i}-1),\,9/q_i\}
       \le c_0:=1/999999 .
       \tag{TD.22}
\]
Definition2.3 and Theorem4.1(2) of \cite{BI:DHP}, pp.~137,141,
apply to all nontrivial irreducible complex characters outside
their scalar loci.  The theorem has no coprime-representation
hypothesis: the coprime hypotheses belong to other results of
that paper.  For a nontrivial representation of a quasisimple
group its scalar locus is its center, since its nonabelian simple
quotient has no proper nontrivial normal subgroup.
Thus (TD.22) applies with the stated exact scope.

Here is an elementary degree consequence which avoids any table
of minimal representation degrees or group orders.
For a nonabelian simple group \(S_i\), let \(\psi\ne1\) be an
ordinary irreducible character of degree \(d\).
If \(dc_0<1\), every value \(\psi(x)\), \(x\ne1\), has all its
Galois conjugates of modulus less than one: those conjugates are
again irreducible characters of degree \(d\), to which (TD.22)
applies.  Character values are algebraic integers, so their
integer norms force \(\psi(x)=0\).
Orthogonality would then imply \(d^2=|S_i|\), contrary to
\(\sum_{\eta\in{\rm Irr}(S_i)}\eta(1)^2=|S_i|\), which also
contains the trivial character.  Therefore
\[
                              d\ge1/c_0=999999 .
                              \tag{TD.23}
\]
The algebraic-integrality, Galois-conjugacy and orthogonality
inputs here are the standard finite-group character facts in
\cite[\S\S2.3--2.4,6.4--6.5,12.1]{OG:Serre1977}.

Now \(S=N/Z(N)=\prod_iS_i\).
Every nontrivial irreducible ordinary \(S\)-module is a tensor
product with at least one nontrivial factor, so (TD.23) gives
the same lower degree \(b=999999\) for \(S\).
We have not asserted a uniform nonidentity trace bound for every
\(S\)-character: that assertion would be false for characters
trivial on a direct factor.
In contrast, the \emph{actual} irreducible \(N\)-module \(W\)
is a tensor product of nontrivial irreducible \(E_i\)-modules.
Faithfulness excludes a trivial \(E_i\)-factor.  Hence every
noncentral element of \(N\) has at least one noncentral component,
and the product character satisfies (TD.17) with \(c=c_0\).

Conjugation on \(\operatorname{End}(W)\) factors through \(S\).
Its invariants have dimension one by irreducibility, and its
nontrivial part has dimension \(a^2-1\).
This part is nonzero, since otherwise every \(N\)-matrix would
be scalar.  Equation (TD.23) consequently gives
\[
                         a^2-1\ge999999,\qquad a\ge1000 .
                         \tag{TD.24}
\]
This argument also covers genuine projective simple-group modules:
no assumption that the center acts trivially on \(W\) has entered.

Apply (TD.18) to every nonscalar normalizer operator.
With \(b=1/c_0\) it gives the exact all-power bound
\[
 \frac{|\operatorname{tr}T|^2}{a^2}
 \le c_0+\frac{1-c_0}{1+1/c_0}
     =\frac{2c_0}{1+c_0}
     =\frac1{500000}<\frac1{170^2}.
 \tag{TD.25}
\]
The same bound holds for every projectively nonidentity power,
because every such power is again a nonscalar normalizer operator.
There is no assumption on the cyclic order or on the outer coset.
An operator which only projectively normalizes the represented
core also normalizes \(N\) itself: the full scalar preimage is
\(\mathbf C^\times N\), whose derived group is \(N\) because
\(N\) is perfect.
If only a finite projective factor image is initially given, take
its inverse image in \(\operatorname{SL}(W)\).  This is a finite
central lift with kernel \(\mu_a\).  It contains \(N\), since
perfectness makes every determinant on \(N\) trivial, and it
normalizes \(N\) by the preceding derived-group argument.
Averaging over this finite lift supplies a common positive
Hermitian form, so all the unitary normalizations used above
are legitimate without a field-of-definition assumption.

Let \(P\) be the actual projective image on \(W\).
It contains the faithful copy of \(S\), and hence is nonsolvable.
Moreover (TD.23) and the ordinary degree-square identity give
\(|S_i|\ge1+999999^2>60\), so \(|P|>60\).
All hypotheses of D-TD02 hold with \(c=1/170\).
It follows that \(e\le504\), whereas
\(e\ge a\ge1000\) by (TD.24), a contradiction.
\end{proof}

\paragraph{What remains.}
D-TD03 excludes cores all of whose defining fields are at least
\(10^{12}\); it does not silently cover a product mixing a small-field
factor with a large-field factor.
The normal core must act irreducibly on the chosen tensor factor.
D-TD01 applies whenever independent arguments bound every projective
factor image, including any future small-degree catalogue.
Neither theorem makes the original inducing field equal to \(K\),
decides a proper Hodge-locus descent, or bounds the arbitrary original
overgroups of D-OG.
